\documentclass[main.tex]{subfiles}


\begin{document}

%from pagenumber to up
% \phantomsection 
% \hypertarget{toc}{}

 

 

% номер секции вручную
\setcounter{section}{44
}
\addtocounter{section}{-1} 

\section{Пар}


Все жидкости\footnote{А также твердые тела.} имеют способность улетучиваться. Например, количество воды в открытом стакане с течением времени, как известно, уменьшается. Частицы жидкости не исчезают бесследно~--- они образуют невидимый пар.

\textbf{Пар}~--- это газ, образованный частицами жидкости, покинувшими ее. В примере с открытым стаканом с водой вылететь из жидкости с ее \emph{свободной поверхности} 
% (то есть с границы раздела жидкости и воздуха)
могут некоторые частицы, имеющие достаточную скорость, чтобы преодолеть силы притяжения к соседним частицам (жидкость испаряется).

\textbf{Испарение}~--- это процесс превращения жидкости в пар, происходящий со свободной поверхности жидкости. \emph{Испарение любой жидкости происходит постоянно при любой температуре}.

Частица пара через некоторое время может вернуться обратно в жидкость~--- пар конденсируется. \textbf{Конденсация}~--- это процесс превращения пара в жидкость. Конденсация пара есть процесс, обратный испарению жидкости.

% Выпадение росы~--- пример конденсации водяного пара. 

Пусть имеются два закрытых сосуда с водой и ее паром (рис.\,\ref{fig:p73}).

    \begin{figure}[H]
    \centering

    \begin{tikzpicture}[font=\small,            line cap=round,                             >={Stealth[inset=0pt,round,length=3mm, angle'=20]},vec/.style={>={Stealth[inset=0pt,round,length=3.5mm, angle'=20]}}]


% \filldraw[lightgray,semitransparent] (0,0) rectangle++ (11,-.3);  
\draw (0,0) --++ (6,0) --++ (0,2.5) --++ (-6,0) -- cycle;

%liquid
\def\dw{.5}
\fill[NavyBlue,nearly transparent] (0,0) ++ (0,\dw) --++ (0,-\dw) --++ (6,0) --++ (0,\dw) -- cycle;


%mol
\begin{scope}[transparency group,semitransparent]    
\def\di{.08}
\foreach \i in {1,...,30}
  \fill[NavyBlue,shift={({2*\di},{\dw+2*\di})}] ({rnd*(6cm-4*\di cm)}, {rnd*(2cm-4*\di cm)}) circle (\di);

% \fill[NavyBlue,shift={({2*\di},{\dw+2*\di})}] ({(6cm-4*\di cm)}, {(1.5cm-4*\di cm)}) circle (\di);

\begin{scope}[yshift=-0.5cm]
    
%evap
\fill[NavyBlue] (1.5,2) circle (\di) coordinate (e1);
\fill[NavyBlue] (3.1,2.1) circle (\di) coordinate (e2);
\fill[NavyBlue] (4.7,2.3) circle (\di) coordinate (e3);


%cond
\fill[NavyBlue] (.5,1.7) circle (\di) coordinate (c1);
\fill[NavyBlue] (2.3,2.05) circle (\di) coordinate (c2);
\fill[NavyBlue] (5.4,1.9) circle (\di) coordinate (c3);
\end{scope}

\end{scope}

%%%arrows
\draw[->,red,shorten <=.08cm,semithick,vec] (c1) -- (.8,.5);
\draw[->,red,shorten <=.08cm,semithick,vec] (c2) -- (2,.5);
\draw[->,red,shorten <=.08cm,semithick,vec] (c3) -- (5.3,.5);

\draw[->,Green,shorten >=.08cm,semithick,vec] (1.1,.5) -- (e1) ;
\draw[->,Green,shorten >=.08cm,semithick,vec] (3.4,.5) -- (e2);
\draw[->,Green,shorten >=.08cm,semithick,vec] (4,.5) -- (e3);

%%%%%%%%%%%%%%%%%%%%%%%%%%%%%%%%%%%%%%
%%%%%%%%%%%%%%%%%%%%%%%%%%%%%%%%%%%%%%

\begin{scope}[xshift=7cm]
% \filldraw[lightgray,semitransparent] (0,0) rectangle++ (11,-.3);  
\draw (0,0) --++ (6,0) --++ (0,2.5) --++ (-6,0) -- cycle;

%liquid
\def\dw{.5}
\fill[NavyBlue,nearly transparent] (0,0) ++ (0,\dw) --++ (0,-\dw) --++ (6,0) --++ (0,\dw) -- cycle;


%mol
\begin{scope}[transparency group,semitransparent]    
\def\di{.08}
\foreach \i in {1,...,15}
  \fill[NavyBlue,shift={({2*\di},{\dw+2*\di})}] ({rnd*(6cm-4*\di cm)}, {rnd*(2cm-4*\di cm)}) circle (\di);

% \fill[NavyBlue,shift={({2*\di},{\dw+2*\di})}] ({(6cm-4*\di cm)}, {(1.5cm-4*\di cm)}) circle (\di);

\begin{scope}[yshift=-0.5cm]
    
%evap
\fill[NavyBlue] (1.5,1.7) circle (\di) coordinate (e1);
\fill[NavyBlue] (3.1,1.8) circle (\di) coordinate (e2);
\fill[NavyBlue] (4.7,2.3) circle (\di) coordinate (e3);


%cond
\fill[NavyBlue] (.5,1.9) circle (\di) coordinate (c1);
\fill[NavyBlue] (2.3,2.05) circle (\di) coordinate (c2);
\fill[NavyBlue] (5.4,1.9) circle (\di) coordinate (c3);
\end{scope}

\end{scope}

%%%arrows
\draw[->,red,shorten <=.08cm,semithick,vec] (c1) -- (.4,.5);
\draw[->,red,shorten <=.08cm,semithick,vec] (c2) -- (2.5,.5);
% \draw[->,red,shorten <=.08cm,semithick,vec] (c3) -- (5.3,.5);

\draw[->,Green,shorten >=.08cm,semithick,vec] (1.1,.5) -- (e1) ;
\draw[->,Green,shorten >=.08cm,semithick,vec] (3.8,.5) -- (e2);
\draw[->,Green,shorten >=.08cm,semithick,vec] (5,.5) -- (e3);

   
\end{scope}



    \end{tikzpicture}
\caption{Пары над жидкостью в сосудах}
\label{fig:p73}
\end{figure}

Зеленые стрелки показывают процессы покидания частицей жидкости; красные стрелки показывают обратные процессы. Соответственно примерам, проиллюстрированным рисунком, можно выделить два вида паров.

\begin{enumerate}
    \item \textbf{Насыщенный} пар (рис.\,\ref{fig:p73}, слева) находится в \emph{динамическом равновесии} со своей жидкостью: число частиц вылетающих из жидкости \emph{равно} числу возвращающихся в нее за то же время. Пар \emph{не может} вместить в себя больше частиц (как бы максимально <<наполнен>> ими).

    Например, насыщенным будет пар над жидкостью, которая долгое время находится в закрытом сосуде.

    \item В системе из \textbf{ненасыщенного} пара (рис.\,\ref{fig:p73}, справа) и его жидкости испарение преобладает над конденсацией: число частиц вылетающих из жидкости \emph{превышает} числу возвращающихся в нее за то же время. 
    Такой пар \emph{может} вместить в себя еще больше частиц.

    К примеру, ненасыщенным можно считать пар над жидкостью, которую только что налили в сухой сосуд и закрыли.
\end{enumerate}

Давление~($P_\text{н.\,п}$) и плотность~($\rho_\text{н.\,п}$) насыщенного пара, например, воды~--- это максимальные давление и плотность, которые может иметь водяной пар при данной температуре. Эти параметры однозначно определяются \emph{только} температурой; они сведены в справочные таблицы для насыщенных паров. Учитывая сказанное, состояние пара можно приближенно описывать уравнением Менделеева"--~Клапейрона.
 
% При \emph{фиксированной температуре} состояния насыщенных и ненасыщенных паров можно приближенно описывать уравнением Менделеева"--~Клапейрона. 
























 
\end{document}